%%%%%%%%%%%%%%%%%%%%%%%%%%%%%%%%%%%%%%%%%
% Plantilla para reporte en LaTeX
% Por Daniel Martínez
% dag.mtz.97@gmail.com
%
%
% Basado en:
%
% WikiBooks (http://en.wikibooks.org/wiki/LaTeX/Title_Creation)
%
% License:
% CC BY-NC-SA 3.0 (http://creativecommons.org/licenses/by-nc-sa/3.0/)
%%%%%%%%%%%%%%%%%%%%%%%%%%%%%%%%%%%%%%%%%

% Declarar el tipo de documento
\documentclass[11pt,letterpaper]{article} 

% Paquetes para idioma, carácteres y símbolos
\usepackage[condensed,math]{iwona}
\usepackage[T1]{fontenc}
\usepackage[spanish, mexico]{babel}
\usepackage{amsfonts}
\usepackage{amssymb}

% Tipo de letra 
% \usepackage{tgbonum}
% \usepackage{tgadventor}

% Márgenes, figuras, gráficos, tablas, listas
\usepackage[top=1.2in, bottom=1.2in, left=1.5in, right=1.5in]{geometry}
\usepackage[justification=centering]{caption}
\usepackage{subcaption}
\usepackage{graphicx}
\usepackage{epstopdf}
\usepackage{tabularx}
\usepackage{float}
\usepackage{enumitem}
%\epstopdfsetup{update} % only regenerate pdf files when eps file is newer
% \usepackage{multirow}
% \usepackage{textcomp}

%\usepackage{mathtools}
%\usepackage{siunitx}
%\usepackage{mathrsfs}

% Otros

% \usepackage{hyperref}
% \usepackage[os=win]{menukeys}
% \usepackage{biblatex}
% \addbibresource{sample.bib}

% \usepackage{fancyhdr}
% \usepackage{listingsutf8}
% \usepackage{courier}
% \usepackage{color} %red, green, blue, yellow, cyan, magenta, black, white

% Configuraciones 
% \input{arduinoLanguage.tex}

% \pagestyle{fancy}

% \renewcommand\lstlistingname{Bloque de código}
% \renewcommand\lstlistlistingname{Índice de códigos}

\setlength{\headheight}{16pt}

% \rhead{\leftmark}
% \lhead{\rightmark}
% \cfoot{\thepage}
% \renewcommand{\headrulewidth}{0.4pt}
% \renewcommand{\footrulewidth}{0.4pt}

\newcommand{\HRule}{\rule{\linewidth}{0.5mm}} % Comando para líneas horizontales

% \definecolor{mygreen}{RGB}{28,172,0} % color values Red, Green, Blue
% \definecolor{mylilas}{RGB}{170,55,241}

%%%%%%%%%%%%%%%%%%%%%%%%%%%%%%%%%%%%%%%%%
% Inicio del documento
%%%%%%%%%%%%%%%%%%%%%%%%%%%%%%%%%%%%%%%%%

\begin{document}
% \hypersetup{pageanchor=false}

% Portada
\begin{titlepage}
    \center
    
    \textsc{\LARGE Universidad de Guadalajara}\\[0.2cm] % Nombre de la universidad
    \textsc{Centro Universitario de Ciencias Exactas e Ingenierías}\\[1cm]
    
    \HRule \\[0.3cm]
    { 
        \LARGE \textbf{Tarea 1} \\[0.3cm]
        \large \textbf{Conceptos sobre sistemas embebidos y microcontroladores} \\[0.3cm]
    }
    \HRule \\[2cm]
    
    \includegraphics[width=6cm]{logo_udg_color.eps}\\[1.25cm] % Logo, necesita paquete graphicx
    
    \center
    \vspace{1cm}
    \emph{Imparte:} \\ % Nombre del profesor
    \textsc{Daniel Giovanni Martínez Sandoval} 
    
    \vfill
    \textbf{I7266 - Programación de Sistemas Embebidos - D10} \\
    Hecho en \textbf{\LaTeX}	
\end{titlepage}

% Índices

% \tableofcontents	% Índice general
% \listoffigures		% de figuras
% \lstlistoflistings  % de listings
% \listoftables		% de tablas

% \hypersetup{pageanchor=true}
% \newpage

% Contenido

\section*{Indicaciones}
\begin{enumerate}[label=(\roman*)]
    \item Realizar las actividades mencionadas a continuación y subirlas en un documento en formato PDF en el apartado de Classroom correspondiente. 
\end{enumerate}
\subsection*{Actividades}
\begin{enumerate}[label=\arabic*)]
    \item Realizar un breve escrito sobre los siguintes temas con base en lo visto en clase:
    \begin{itemize}
        \item Qué es un sistema embebido
        \item Qué es un microcontrolador
        \item A qué nos referimos con \textit{arquitectura}
    \end{itemize}
    \item Investiga sobre 5 diferentes arquitecturas populares o actualmente usadas, escribe un párrafo con lo más notable de cada una.
    \item Lista 10 sistemas embebidos que uses en tu día a día.
\end{enumerate}
% \subsection*{Generalidades}
% \begin{itemize}
%     \item Las actividades escritas se entregarán en un documento en formato PDF.
%     \item La evidencia de prácticas se entregará también en un documento en formato PDF, además de un archivo comprimido ZIP que incluya el código fuente, código objeto e imágenes como evidencia de la actividad realizada.
%     \item En trabajos de investigación es obligatorio incluir las fuentes, el formato es libre pero se recomienda APA.
% \end{itemize}

% \begin{thebibliography}{10}

%     \bibitem{1} Georg Goubau, \textit{Multi-element monopole antennas}, Wiesbaden, Dieterich'sche Verlagsbuchhandlung, 1976.
    
%     \bibitem{2} C.B. Ravipati, S.R. Best, \textit{The Goubau Multi Element Monopole Antenna - Revisited}, Hampton, Richmond upon Thames, 2007.
    
%     \bibitem{3} Ivor L. Morrow, Ryan Walker, Mark Finnis, William Whittow, \textit{Investigation of a Disk-Loaded Monopole Antenna with Extended Bandwidth Matching}, Addison Wesley, Massachusetts, 1994.
    
%     \bibitem{4} Constantine A. Balanis, \textit{Antenna Theory - Analysis and Design}, Arizona State University, 3rd edition, pp. 611–650, 1986.
    
%     \bibitem{5} R. C. Hansen, \textit{Electrically Small, Superdirective and Superconducting Antennas}, Wiley-Interscience, 2006.
  
% \end{thebibliography}
\end{document}

% Esta es una plantilla para poner dos imágenes lado a lado en una sola figura
% \begin{figure}
%     \centering
%     \begin{subfigure}[b]{0.45\textwidth}
% 			\centering
% 			\captionsetup{justification=centering}
% 			\includegraphics[width=0.7\textwidth]{example-image} % Imagen
% 			\caption{Pie de imagen 1} % Pie de imagen
% 			\label{img:1} % Etiqueta
%     \end{subfigure}
%     ~ %add desired spacing between images, e. g. ~, \quad, \qquad, \hfill etc.
%       %(or a blank line to force the subfigure onto a new line)
%     \begin{subfigure}[b]{0.45\textwidth}
% 			\centering
% 			\captionsetup{justification=centering}
% 			\includegraphics[width=0.7\textwidth]{example-image} % Imagen
% 			\caption{Pie de imagen 2} % Pie de imagen
% 			\label{img:2} % Etiqueta
%     \end{subfigure}
%     \caption{Este es un ejemplo de figura, con subfiguras.}\label{fig:1} % Pie de figura y etiqueta
% \end{figure}

% \paragraph{Elaboración del prototipo} Cuando las pruebas y los resultados nos sean satisfactorios y además estos mismos se relacionen con los propuestos en los documentos comenzaremos a hacer simulaciones mas a fondo con la información que nosotros mismos vallamos obteniendo.

% \paragraph{Simulación y resultados} Los resultados de las simulación deberán ser lo mas parecidos a los de los documentos, además intentaremos mejorar estos mismos con lo aprendido en el curso y nuestros conocimientos adquiridos.

% \paragraph{Antena prototipo} Ya que nuestro primer prototipo coincida satisfactoriamente con los documentos además de con los cálculos hechos, tendremos que hacer un modelo definitivo de la antena y este sera el que presentara todos los resultados y mediciones, además de que este sera el que tenga un mejor rendimiento en todos sus parámetros y en base a este prototipo daremos nuestras colusiones finales.


% \paragraph{Documentación del trabajo} El ultimo paso a realizar sera la documentación. Conforme vayamos avanzando en la construcción de la antena, se tomara nota de los resultados tanto satisfactorios como erróneos, y de esta forma tener una mejor visualización de como fue avanzando el proyecto y la construcción.

